\section{Основные определения}
\label{sec:FLT_BasicDefinitions}
\tikzsetfigurename{FLT_BasicDefs_}

\begin{definition}[Алфавит]
    \emph{Алфавит}~--- это конечное множество.
    Элементы этого множества будем называть \emph{символами}.
\end{definition}

\begin{example}
    Примеры алфавитов.
    \begin{itemize}
        \item Латинский алфавит $\Sigma = \{a, b, c, \dots, z\}$.
        \item Кириллический алфавит $\Sigma = \{\text{а, б, в, \dots, я}\}$.
        \item Алфавит натуральных чисел в шестнадцатеричной записи
              \[\Sigma = \{0, 1, 2, 3, 4, 5, 6, 7, 8, 9, A, B, C, D, E, F\}.\]
    \end{itemize}
\end{example}

Традиционное обозначение для алфавита~--- $\Sigma$.
Также мы будем использовать различные прописные буквы латинского алфавита.
Для обозначения символов алфавита будем использовать строчные буквы латинского алфавита: $a, b, \dots, x, y, z$.

Будем отождествлять символы и слова длины один.
Тогда можно считать, что всегда определена операция конкатенации $\cdot: \Sigma^* \times \Sigma^* \to \Sigma^*$.
Здесь $\Sigma^*$~--- замыкание множества%
\sidenote{
Именно здесь нам важно сичтать, что символ и слово длины 1~--- это одно и то же. Иначе для замыкания и для <<привычной>> конкатенации пришлось бы задавать отдельные операции и выражать одну через другую.
}%
$\Sigma$ относительно операции $\cdot$ (см. ~\ref{def:set_closure}).
При записи выражений символ точки (обозначение операции конкатенации) часто будем опускать: $a \cdot b = ab$.

\begin{definition}[Слово]
    \emph{Слово} над алфавитом $\Sigma$~--- это конечная конкатенация символов алфавита $\Sigma$: $\omega = a_0 \cdot a_1 \cdot \dots \cdot a_m$, где $\omega$~--- слово, а $a_i \in \Sigma$ для любого $i$.
    Для обозначения пустого слова (слова, содержащего ноль символов) будем использовать специальный символ $\varepsilon, \varepsilon \notin \Sigma$.
\end{definition}

\begin{definition}[Длина слова]
    Пусть $\omega = a_0 \cdot a_1 \cdot \dots \cdot a_m$~--- слово над алфавитом $\Sigma$.
    Будем называть $m + 1$ \emph{длиной слова} и обозначать как $|\omega|$.
    Длина пустого слова равна нулю.
\end{definition}

При работе с линейным входом (строками) нам будет удобно ввести два связанных понятия позиции.

\begin{definition}[Позиция символа]
    Пусть $\omega = a_0 a_1 \dots a_{n-1}$~--- слово длины $n$ над алфавитом $\Sigma$.
    \emph{Позицией символа} $a_i$ в слове $\omega$ будем называть число $i$.
\end{definition}

\begin{definition}[Позиция в строке]
    \label{def:PositionInString}
    Пусть $\omega = a_0 a_1 \dots a_{n-1}$~--- слово длины $n$ над алфавитом $\Sigma$.
    \emph{Позиции в строке} нумеруются от $0$ до $n$ и располагаются между символами:
    \[
    {}_{0} {a_0}_{1} {a_1}_{2} \dots {}_{n-1} {a_{n-1}}_{n}.
    \]
\end{definition}

Нетрудно видеть, что символ $a_i$ находится между позициями $i$ и $i+1$.
Таким образом, позиция символа $a_i$ совпадает с левой из двух определяющих его позиций в строке.
Данное соглашение упрощает единообразное описание процедур анализа строк, в которых состояние анализатора изменяется при переходе от одной позиции к другой, сопровождаемом чтением символа.

\begin{definition}[Количество вхождений символа]
    Пусть $\omega$~--- слово над алфавитом $\Sigma$. Пусть $t \in \Sigma$~--- некоторый символ из алфавита.
    Тогда будем обозначать количество вхождений символа $t$ в слове $\omega$ как $|\omega|_t$.
\end{definition}

\begin{definition}[Язык]
    \emph{Язык} над алфавитом $\Sigma$~--- это множество слов над алфавитом $\Sigma$: $L_{\Sigma} \subseteq \Sigma^{*}$.
\end{definition}

Иначе говоря, любой язык над алфавитом $\Sigma$ является подмножеством универсального множества $\Sigma^*$~--- множества всех слов над алфавитом $\Sigma$.

\begin{example}
    Примеры языков.

    \begin{itemize}
        \item Язык целых чисел в двоичной записи
              \[\{0, 1, 10, 11, 100, 101, \dots\}.\]
        \item Язык всех правильных скобочных последовательностей. Он же~--- частный случай \textbf{языка Дика}.
              \[\{(), (()), ()(), (())(), \dots\}.\]
        \item \textbf{Язык Дика} на двух типах скобок.
              \[\{(), (()[]), [](), (()[])(), [], \dots\}.\]
    \end{itemize}
\end{example}


\mytodo{Добавить ссылки на работы в sidenote-ы ниже.}
Заметим, что язык не обязан быть конечным множеством, в то время как алфавит в нашей области всегда конечен%
\sidenote{Существуют ситуации, когда возникают бесконечные алфавиты.}
и изучаем мы конечные слова%
\sidenote{
    Существуют ситуации, когда возникают бесконечные слова.
    Например работы по обработке потоков.
}.

Можно выделить следующие основные \emph{способы задания языков}.
\begin{itemize}
    \item Перечислить все элементы.
          Такой способ работает только для конечных языков.
          Перечислить бесконечное множество за конечное время не получится.
    \item Задать генератор~--- процедуру, которая на каждый свой вызов возвращает очередное слово языка.
    \item Задать распознаватель~--- процедуру, которая по данному слову может определить, принадлежит оно заданному языку или нет.
\end{itemize}

Стоит отметить, что существуют и другие способы задания.
Например, язык может определяться как решение некоторого \emph{языкового уравнения}~\sidecite{Leiss1999, OKHOTIN2025105344}.
